% M6 L1マトリクス掃引: 代表クラスオプション一式（非既定値）のコンパイル検証
% 既知制限（pdfTeXバックエンド）: 値にリテラルzw/zhを含む\jlreqsetupキー
% （footnote_indent, footnote_second_indent, quote_indent等）は，キー値を
% 後段の\dimexprで評価する実装のためpdfTeXでは「Illegal unit」になる．
% これらはバックエンド側（jlreq-pdftex）の拡張が必要なためM6では
% 「limitation」として文書化する（docs/pdftex-compatibility.md参照）．
\documentclass[pdflatex,paper=b5,fontsize=12pt,line_length=40zw,number_of_lines=10,head_space=15mm,foot_space=20mm,gutter=15mm,fore_edge=20mm,baselineskip=20pt,column_gap=2zw,twoside]{jlreq}
\jlreqsetup{
  caption_font=\small,
  caption_label_format={表#1:},
  itemization_itemsep=0.2\baselineskip,
  warichu_opening=（, warichu_closing=）,
}
\begin{document}
\section{版面オプションの掃引}
紙・行長・行数・天地・のど・小口・行送り・段間を非既定値にした文書．
割注\warichu{割注一行目\\割注二行目}を含む．
\begin{itemize}\item 箇条書きの項目間空き検証．\end{itemize}
\begin{table}[htbp]\centering\caption{キャプション書式検証}
\begin{tabular}{cc} あ & い \\ う & え \\\end{tabular}\end{table}
\subsection{小見出し}
\end{document}
